Reported accuracies for automatic cough detection sit in the high nineties, yet audits find
such models are rarely validated outside their development corpus and often separated by a
confound. Which factor is responsible has not been isolated, because label provenance,
negative class difficulty and seed variation ordinarily move together. This study fixes the
preprocessing, architecture and training protocol, and varies only the provenance of the
label. Three corpora are used, labelled respectively by a participant's instruction, by a
human listener against an audio ontology, and by another classifier's score. The class
prior is matched in all three, partitions are speaker disjoint wherever an identifier
permits, every condition runs under five seeds, and conditions are compared by a paired
test on per recording decisions. Within the prompted corpus a residual network reaches 98.6
per cent balanced accuracy on held out speakers. The same weights, applied without
retraining to human annotated audio at the same prior, reach 70.2 per cent: specificity
falls to 84.3 per cent, sensitivity to 56.0. The loss is one of discrimination rather than
calibration, the area under the curve falling from 0.998 to 0.749, so moving the threshold
recovers little of it. Sneezes are called coughs more often than coughs are, locating the
failure in a negative class holding no forced expirations. Wavelet shrinkage does not help,
removing a median of 0.67 per cent of signal power. The external set is small and single
source, so the fall is better established in direction than in size.
